\documentclass[spanish,letterpaper,oneside]{article}

% Datos del documento
\def\documenttitle {Metodologías de la NEM y decisiones de enseñanza}
\def\documentsubtitle {Fundamentos para la enseñanza y el aprendizaje I}
\def\documentsubject {Actividad 7 · Cuadros comparativos y propuesta de aplicación}
\def\documentauthor {Mtro. Martín Jonathan de la Cruz Muñoz}
\def\coursename {Fundamentos para la enseñanza y el aprendizaje I}
\def\coursecode {MAEM}

% Identidad institucional
\def\universityname {Instituto de Investigación, Innovación y Estudios de Posgrado para la Educación del Estado de Nuevo León}
\def\universityfaculty {Maestría en Aprendizaje y Enseñanza de las Matemáticas}
\def\universitydepartment {Fundamentos para la enseñanza y el aprendizaje I}
\def\universitydepartmentimage {img/iiiepe/marca-iiiepe.png}
\def\universitydepartmentimagecfg {height=1.75cm}
\def\universitylocation {Monterrey, Nuevo León}
\def\studentfullname {Mtro. Martín Jonathan de la Cruz Muñoz}
\def\studentid {00874}
\def\studentemail {martin.delacruz@campus.iiiepe.edu.mx}
\def\programperiod {1er tetramestre, marzo--julio de 2026}
\def\programmodality {Modalidad presencial}
\def\programshort {MAEM}

\def\authortable {
  \begin{tabular}{ll}
    \textbf{Alumno:} & \studentfullname \\
    Matrícula: & \studentid \\
    Correo institucional: & \studentemail \\
    Unidad: & \coursename \\
    Programa: & \programshort \\
    Periodo: & \programperiod \\
    & \\
    \multicolumn{2}{l}{Fecha de realización: 13 de octubre de 2026} \\
    \multicolumn{2}{l}{\universitylocation}
  \end{tabular}
}

\input{template}

\usepackage{helvet}
\renewcommand{\familydefault}{\sfdefault}
\usepackage{array}
\usepackage{longtable}
\usepackage{pdflscape}
\usepackage{tikz}
\usepackage[most]{tcolorbox}
\usetikzlibrary{arrows.meta,positioning,calc,fit,backgrounds}

\definecolor{iiiepeNavy}{HTML}{163B5C}
\definecolor{iiiepeBlue}{HTML}{1F6F9F}
\definecolor{iiiepeAqua}{HTML}{20B7A8}
\definecolor{iiiepeOrange}{HTML}{F28C28}
\definecolor{iiiepePaper}{HTML}{F6F8F7}
\definecolor{iiiepeInk}{HTML}{1E2A32}
\def\portraittitlecolor {iiiepeNavy}
\setcitestyle{authoryear,open={(},close={)}}

\newcommand{\insertcoverwatermark}{%
  \AddToShipoutPictureBG*{%
    \begin{tikzpicture}[remember picture,overlay]
      \node[opacity=.09,inner sep=0pt] at (current page.center)
      {\includegraphics[width=.58\paperwidth]{img/iiiepe/logo-puntos.png}};
    \end{tikzpicture}%
  }%
}

\newtcolorbox{reportelecturabox}[1][]{%
  enhanced,breakable,
  colback=iiiepePaper,colframe=iiiepeNavy,
  boxrule=1pt,leftrule=4pt,arc=2pt,
  left=8pt,right=8pt,top=6pt,bottom=6pt,
  before skip=6pt,after skip=10pt,
  fonttitle=\bfseries\color{white},coltitle=white,
  colbacktitle=iiiepeNavy,
  attach boxed title to top left={xshift=8pt,yshift=-2pt},
  boxed title style={colback=iiiepeNavy,colframe=iiiepeNavy,arc=1pt},
  title={Reporte de lectura},#1}

\newtcolorbox{fichalecturabox}[1][]{%
  enhanced,breakable,
  colback=iiiepeAqua!7,colframe=iiiepeAqua!70!iiiepeNavy,
  boxrule=.6pt,arc=1pt,
  left=7pt,right=7pt,top=4pt,bottom=4pt,
  fontupper=\small,before upper={\setstretch{1.0}},
  fonttitle=\bfseries\color{iiiepeNavy},coltitle=iiiepeNavy,
  colbacktitle=iiiepeAqua!7,
  title={Ficha de la lectura},#1}



\usepackage{tabularx,colortbl,pgfplots,enumitem,needspace}\pgfplotsset{compat=1.18}\definecolor{iiiepePale}{HTML}{EEF5F8}
\setlength{\parskip}{5pt}\setlength{\parindent}{0pt}
\begin{document}\insertcoverwatermark\templatePortrait\templatePagecfg\templateFinalcfg\fancyfoot{}\fancyfoot[L]{\scriptsize IIIEPE · MAEM}\fancyfoot[C]{\scriptsize\thepage}\fancyfoot[R]{\scriptsize Actividad 7}\setstretch{1.3}\raggedbottom
\section{Elegir una metodología según el problema}
Las sugerencias de la SEP son alternativas para organizar proyectos educativos y no un recetario obligatorio. La elección debe atender al propósito, al problema y a las condiciones de participación. Se distinguen proyectos comunitarios, indagación con STEAM, aprendizaje basado en problemas y aprendizaje servicio. La sigla ABP es ambigua: en este reporte se escribe el nombre completo para diferenciar proyectos y problemas.
\section{Cuadro de cuatro entradas}


{\small\setstretch{1.1}\renewcommand{\arraystretch}{1.25}
\begin{longtable}{|>{\raggedright\arraybackslash}p{\dimexpr0.3\linewidth-2\tabcolsep-2\arrayrulewidth\relax}|>{\raggedright\arraybackslash}p{\dimexpr0.35\linewidth-2\tabcolsep-2\arrayrulewidth\relax}|>{\raggedright\arraybackslash}p{\dimexpr0.35\linewidth-2\tabcolsep-2\arrayrulewidth\relax}|}\hline
\rowcolor{iiiepePale} \textbf{Metodología} & \textbf{Cómo se trabaja} & \textbf{Palabras clave} \\\hline
\endfirsthead
\hline
\rowcolor{iiiepePale} \textbf{Metodología} & \textbf{Cómo se trabaja} & \textbf{Palabras clave} \\\hline
\endhead
Proyectos comunitarios & Parte de una situación del entorno; se acuerda un plan, se produce una respuesta y se comunica y revisa con la comunidad. La SEP organiza planeación, acción e intervención con momentos de seguimiento. & Comunidad, producción, contexto, participación, comunicación. \\\hline
Indagación con STEAM & Se formula una pregunta, se diseña una indagación, se organizan resultados, se presentan y se reflexiona. Integra ciencias, tecnología, ingeniería, artes y matemáticas cuando contribuyen al problema. & Pregunta, hipótesis, evidencia, modelación, diseño, metacognición. \\\hline
\end{longtable}}



{\small\setstretch{1.1}\renewcommand{\arraystretch}{1.25}
\begin{longtable}{|>{\raggedright\arraybackslash}p{\dimexpr0.3\linewidth-2\tabcolsep-2\arrayrulewidth\relax}|>{\raggedright\arraybackslash}p{\dimexpr0.35\linewidth-2\tabcolsep-2\arrayrulewidth\relax}|>{\raggedright\arraybackslash}p{\dimexpr0.35\linewidth-2\tabcolsep-2\arrayrulewidth\relax}|}\hline
\rowcolor{iiiepePale} \textbf{Metodología} & \textbf{Cómo se trabaja} & \textbf{Palabras clave} \\\hline
\endfirsthead
\hline
\rowcolor{iiiepePale} \textbf{Metodología} & \textbf{Cómo se trabaja} & \textbf{Palabras clave} \\\hline
\endhead
Aprendizaje basado en problemas & Se presenta una situación, se recuperan saberes, se define el problema, se organiza la experiencia, se investiga y se analizan resultados. El proceso exige construir explicaciones y decisiones fundamentadas. & Problematización, saberes previos, investigación, argumentación, solución. \\\hline
Aprendizaje servicio & Se reconoce una necesidad comunitaria y se acuerda un servicio ligado a aprendizajes. Se organiza, realiza, evalúa y comparte la experiencia. El servicio responde a una necesidad y no se reduce a una actividad asistencial. & Reciprocidad, necesidad, servicio, aprendizaje, responsabilidad, reflexión. \\\hline
\end{longtable}}


\section{Características y roles docentes y estudiantiles}


{\small\setstretch{1.1}\renewcommand{\arraystretch}{1.25}
\begin{longtable}{|>{\raggedright\arraybackslash}p{\dimexpr0.3\linewidth-2\tabcolsep-2\arrayrulewidth\relax}|>{\raggedright\arraybackslash}p{\dimexpr0.35\linewidth-2\tabcolsep-2\arrayrulewidth\relax}|>{\raggedright\arraybackslash}p{\dimexpr0.35\linewidth-2\tabcolsep-2\arrayrulewidth\relax}|}\hline
\rowcolor{iiiepePale} \textbf{Metodología} & \textbf{Papel docente} & \textbf{Papel del estudiante} \\\hline
\endfirsthead
\hline
\rowcolor{iiiepePale} \textbf{Metodología} & \textbf{Papel docente} & \textbf{Papel del estudiante} \\\hline
\endhead
Proyectos comunitarios & Ayuda a delimitar una situación pertinente, acompaña el plan, ofrece criterios y facilita diálogo con destinatarios. & Participa en decisiones, investiga, produce y explica una respuesta; revisa su pertinencia con otros. \\\hline
Indagación con STEAM & Apoya preguntas investigables, controla condiciones de seguridad, ayuda a examinar datos y conceptos. & Predice, diseña pruebas, registra, modela, contrasta y comunica; distingue observación de interpretación. \\\hline
\end{longtable}}



{\small\setstretch{1.1}\renewcommand{\arraystretch}{1.25}
\begin{longtable}{|>{\raggedright\arraybackslash}p{\dimexpr0.3\linewidth-2\tabcolsep-2\arrayrulewidth\relax}|>{\raggedright\arraybackslash}p{\dimexpr0.35\linewidth-2\tabcolsep-2\arrayrulewidth\relax}|>{\raggedright\arraybackslash}p{\dimexpr0.35\linewidth-2\tabcolsep-2\arrayrulewidth\relax}|}\hline
\rowcolor{iiiepePale} \textbf{Metodología} & \textbf{Papel docente} & \textbf{Papel del estudiante} \\\hline
\endfirsthead
\hline
\rowcolor{iiiepePale} \textbf{Metodología} & \textbf{Papel docente} & \textbf{Papel del estudiante} \\\hline
\endhead
Aprendizaje basado en problemas & Propone o acompaña un problema relevante, orienta búsqueda y pide justificar alternativas. & Identifica lo que sabe y necesita saber; contrasta fuentes, argumenta y revisa una solución. \\\hline
Aprendizaje servicio & Vincula objetivos curriculares con una necesidad acordada y acompaña evaluación y reciprocidad. & Escucha a destinatarios, participa en el servicio y explica qué aprendió y cómo responde a la necesidad. \\\hline
\end{longtable}}


\section{Ejemplo propuesto para mi contexto}
Propongo indagación con STEAM para estudiar un sistema de medición de bajo costo. En preparatoria puede centrarse en proporcionalidad y gráficas; en ingeniería, en calibración, regresión y análisis del error. La necesidad y la complejidad deben adecuarse al grupo y a los recursos disponibles. La secuencia no presupone que ya exista un problema medido en una institución.
En el inicio se plantea cómo decidir si un instrumento mide de manera confiable. Después se identifican variables y se diseña una prueba segura. En el desarrollo se registran lecturas y referencias, se construyen tabla y gráfica y se revisa el modelo. La comunicación explica resultados y límites; la metacognición identifica un cambio de razonamiento. El diseño integra tecnología solo si aporta a la investigación y ofrece una alternativa con materiales sencillos.
La evaluación considera formulación del problema, trazabilidad de datos, interpretación matemática y calidad del argumento. Se solicita una revisión después de la devolución. La distribución de roles permite observar participación; no se asume que todos aprendieron porque el equipo entregó un prototipo.
\section{Criterios para elegir sin confundir metodologías}
La selección puede formularse como una decisión: si el centro es elaborar una respuesta pública con destinatarios, conviene explorar proyectos comunitarios; si la pregunta exige obtener y contrastar evidencia, indagación; si el obstáculo requiere estudiar alternativas y construir una solución fundamentada, aprendizaje basado en problemas; si existe una necesidad acordada y un servicio recíproco vinculado con contenidos, aprendizaje servicio. Las fronteras pueden relacionarse, pero cada diseño debe explicar su centro.
En el material de la SEP (2022a, 2022b), las propuestas se ofrecen como posibilidades y no como una prescripción que anule la autonomía docente. Por eso justificaría la elección por la actividad intelectual y social que se busca, los recursos y el tiempo. Un mismo problema de consumo de agua podría investigarse con mediciones o convertirse en servicio comunitario; el nombre del tema no decide la metodología.
Para mi contexto, la indagación de un sistema de medición permite organizar preguntas, modelos y contraste. La dimensión humanista se concreta al acordar una pregunta comprensible, distribuir decisiones y ofrecer alternativas de acceso. El cierre debe mostrar qué aprendió cada integrante y cómo se utilizó una devolución para revisar, además de comunicar el producto colectivo.

\section{Comparación y postura}
Los cuatro enfoques comparten participación y contextualización, pero su centro difiere: producción comunitaria, indagación, resolución de un problema o servicio recíproco. Pueden relacionarse, aunque no es necesario combinarlos todos en una secuencia. Mi formación industrial favorece planificar y verificar; el enfoque humanista me exige añadir diálogo sobre necesidades y condiciones de aprendizaje. Elegiría la metodología que permita sostener esa relación y comprobar qué comprenden los estudiantes.

\needspace{5\baselineskip}\section*{Referencias}
{\setstretch{1.15}\setlength{\parskip}{6pt}
\par\hangindent=1.27cm\hangafter=1  Secretaría de Educación Pública. (2022a). \textit{Plan de Estudio para la educación preescolar, primaria y secundaria}. SEP.
\par\hangindent=1.27cm\hangafter=1  Secretaría de Educación Pública. (2022b). \textit{Sugerencias metodológicas para el desarrollo de los proyectos educativos} [Extracto de \textit{Avance del contenido para el libro del docente. Primer grado}, material en proceso de edición, pp. 64--82, 23--24]. SEP.
}

\end{document}
